\newcommand{\alphavalue}{\sqrt{\frac{1}{8}}}
\newcommand{\betavalue}{\sqrt{\frac{7}{8}}}

\section{Quantencomputer}

Quantencomputer sind keine neuen Computer, sie können in der Lage manche Ressourcenlastingen Berechnungen effizienter zu lösen als es ein traditioneller Computer könnte.
Das Jahr in dem die ersten funktionsfähigen Quantencomputer geben wird, wird als als Y2Q bezeichnet, der genaue Tag als Q-Day.\cite{schmehQuantencomputer2026}
Quantencomputer sind noch weit entfernt von dem Zustand in dem sie wirklich eine Gefahr für die Moderne Kryptographie darstellen,
als vorkehrungsmaßnahme gibt es aber bereits algorthmen die theoretisch einem solchen Angriff standhalten sollen.

Ein Quantencomputer macht sich die Gesetze der Quantenmechanik zu nutze, genauer gesagt Superposition und Quantenverschränkung.
Anstelle von Bits werden sogenannte Qubits verwendet.

\subsubsection{Qubits}
Ein Qubit kann, mit einer Wahrscheinlichkeit, Zustände von 0 oder 1 annehmen. Um ein Qubit als Zahl darzustellen, macht man das als komplexe Zahl.
Es gibt die Wahrscheinlichkeit für eine null, der Realteil $\alpha$ und die Wahrscheinlichkeit auf eine eins, der Imaginärteil $\beta=\betavalue$.
Dass ein Qubit Eins oder Null gleichzeitig sein kann, bis man es observiert, ist Superposition. Wenn man es observiert, nimmt es eindeutig einen Wert an.
Man stellt das Qubit als vektor $\psi$ da. Nehmen wir an wir haben $\alpha=\alphavalue$ und $\beta=\betavalue$.

\begin{align}
    \ket{\psi}\ =\ (\sqrt{\frac{1}{8}})\ket{0}\ +\ (\sqrt{\frac{7}{8}})\ket{1}
\end{align}

Das bedeutet wir haben die Wahrscheinlichkeiten:
\begin{align}
    &{\left| \alpha \right|}^2\ =\ \alphavalue^2\ =\ \frac{1}{8}\\
    &{\left| \beta \right|}^2\ =\ \betavalue^2\ =\ \frac{7}{8}
\end{align}

Die Wahrscheinlichkeiten ergeben in der Summe $\left| \alpha \right|^2 + \left| \beta \right|^2\ =\ 1$.

Wir können den Vektor der Wahrscheinlichkeiten in einem Koordinatensystem darstellen.
\newpage

\begin{figure}[htbp]
\centering
\begin{tikzpicture}
% Achsen
\draw[->] (0,0) -- (5,0) node[right] {$\alpha$};
\draw[->] (0,0) -- (0,4) node[above] {$\beta$};

% Linie vom Ursprung aus
\draw[thick, blue] (0,0) -- (3,2);
\end{tikzpicture}
\caption{Qubit in Koordinatensystem darstellung}
\label{fig:koordinatensystem}
\end{figure}

Wenn wir mehr als ein Qubit haben, wie hier z.B. 128, können wir damit $2^{7}=128$ Zustände darstellen.

\begin{align}
    &\alpha_0\ket{0000000}+\alpha_1\ket{0000001}+\dots+\alpha_{127}\ket{1111111}
\end{align}

Alle Wahrscheinlichkeiten müssen zusammen wieder eins ergeben.
\begin{align}
    &{\left|\alpha_0\right|}^2+{\left|\alpha_1\right|}^2+\dots+{\left|\alpha_{127}\right|}^2=1
\end{align}
Man kann sich den Vektor als Punkt im 128-dimensionalem Koordinatensystem darstellen. Die Achse zu der die Linie am ehesten Senkrecht steht, ist die höchte Wahrscheinlichkeit.
Durch den Grover Algorithmus lässt sich eine Annäherung machen.

\subsubsection{Grover's Algorithmus}
Ich weiß nicht was ich hier schreiben soll. Ganz ehrlich\dots